\chapter{知识的演进}
\label{chap:ket-overview}

前面七章都默认“\textbf{有一个\emph{固定}的任务、训练一次}”。但真实系统几乎从不安静：模型上线后，\textbf{世界在变、需求在变、数据在变}。设想你已经做好一个花卉识别系统，能从照片判断物种，并且它已经在一个园区里稳定运行了一段时间。接下来的一年里，接连发生了几件事：\textbf{雨季来了}，照片里的阴影、反光、背景色调都和训练时不一样；\textbf{园区逐月引进新物种}，每隔一阵就得让系统认一种原来没有的花；\textbf{机构负责人由甲交接给乙}，附带的问答服务不能再把现任负责人答错。

面对这些，你\textbf{手里已经有一个训练好、正在服务的旧模型}——它显然是个\textbf{宝贵的现成起点}（里面已经学到了花瓣纹理、轮廓、颜色这些东西），但它\textbf{\emph{不能}直接替你回答三个问题}：
\begin{enumerate}\item \textbf{环境变了以后，旧模型里\emph{哪些判断}仍然成立}？（依赖物体本身的线索或许还在，依赖当地背景、光照的线索可能已经失效）；\item \textbf{逐月加入新物种时，怎么学新东西而\emph{不冲掉}旧能力}？（学会了新花、忘了老花）；\item \textbf{负责人换了，到底该\emph{改哪个对象}}？（改知识库条目、改模型参数、还是只在这次提问里说明），改完又怎么验证\textbf{历史问题}仍答对？\end{enumerate}
\textbf{本部分（第8–11章）要研究的，正是“一个\emph{已经在跑}的系统，怎样在\emph{变化}中\emph{继续}正确、继续可学}”——核心立场是：\textbf{绝不要把每次新变化都当成“从零开始训练”}，而要先问“我\textbf{已经学到的东西，哪些能留、怎么改、改了会不会丢、以后还学不学得动}”。

\textbf{阅读本章你只需要}：任务、模型、损失、训练、泛化的基本概念，以及“神经网络中间层\emph{表示}”的直观认识——\textbf{\emph{不需要}你已经会任何具体的持续学习或知识编辑算法}，那些是后面章节才展开的。本章是\textbf{整个“知识演进”部分的\emph{路线图}}，它按下面的顺序帮你搭好框架：
\begin{itemize}\item \textbf{先在“场景”里\emph{辨认变化发生在哪}}（是数据分布变了、要陆续学新类别、还是一条世界事实过期了）；\item \textbf{再用“基本概念”说清三件事}：\textbf{经验\emph{保存在哪}}（样本？参数？事实？）、\textbf{一次更新\emph{实际改了哪个对象}}、\textbf{系统\emph{最终具备}了什么能力}，以及\textbf{学习算法\emph{能读取}哪些样本/模型/反馈}（权限）；\item 然后“\textbf{研究内容}”把一个模糊的更新请求\textbf{拆成几个\emph{可判断}的决策}（改哪、怎么保旧、怎么验）；\item 最后“\textbf{研究历史}”说明这些决策\textbf{为什么会\emph{分成几条并行的研究线索}}。\end{itemize}
所以\textbf{本章给的是“问题地图”和“概念边界”，\emph{不给}具体算法}——“具体怎么选、怎么更新、怎么验证”是后续第9–11章的事。

这里要先\textbf{划清一条重要边界}，因为本卷紧接着还有一个部分叫“知识的迁移与共享”（第12–15章）。\textbf{“演进”和“迁移”是两种不同起点的问题}：
\begin{itemize}\item \textbf{知识的演进}（本部分）：\textbf{盯住\emph{同一个}持续运行的系统}，问它在\emph{时间}上怎样被更新、又怎样\emph{保持}旧能力和可学习性。关键词是“\textbf{随时间、边用边改、别忘旧}”。\item \textbf{知识的迁移与共享}（下一部分）：问\textbf{一个任务上\emph{已经学到}的经验，怎样拿去帮\emph{另一个}新任务}（迁移），以及\textbf{多个任务怎样\emph{一起}训练、共享同一套表示}（多任务共享）。关键词是“\textbf{跨任务、老帮新、一起学}”。\end{itemize}
两者当然会交织：\textbf{逐月加入新物种}既是“系统在演进”，每一阶段又都在“把旧物种的经验迁移给新物种”。本书的处理是：\textbf{本部分负责“\emph{时间维度上的保持与维护}”（含用迁移当手段时怎样不遗忘），迁移“\emph{跨任务帮助}”的专门机制（来源/目标、冻结微调、领域自适应、元学习、少样本等）统一下一部分讲}。本章在概念上会点到它们的名字，但不展开算法。

和第一部分的衔接也要清楚：\textbf{第~\ref{part:efficient-knowledge-use}部分（第1–7章）解决的是“\emph{知识从哪来、怎么融合、质量如何}”}——答案、关系、约束、训练目标是怎么\emph{取得}的。\textbf{本部分\emph{换了个起点}}：假设这些东西\textbf{\emph{已经学到了}，系统\emph{已经在跑}}，我们从它\textbf{已经形成的“表示、模型参数、预测行为、记忆”出发}，考察它们在世界变化时\textbf{怎样\emph{继续}被维护}。

打个比方：第一部分是“\textbf{怎么招到、培训好一支队伍并做质检}”，本部分是“\textbf{这支队伍已经上岗后，业务随季节变化、还不断来新人新业务，怎么\emph{让老员工学会新活、又不丢掉老业务、也别越改越僵化}}”。

最后声明本部分的\textbf{文献态度}：这是\textbf{教程式综述，选材截止 2026 年 10 月，\emph{不}声称穷尽这一时期所有研究}。三类文献各有分工：\textbf{经典工作}用来\textbf{建立“问题是什么、机制是什么”}；\textbf{新近工作}用来说明\textbf{“使用条件或研究对象\emph{怎样被扩展}”}（方法边界被推到哪）；\textbf{比较研究和\emph{负面}结果}用来\textbf{\emph{限制}“收益”结论}（什么情况下不灵）。\textbf{两条不能越界的读法}：其一，选材\emph{不能}由“论文数量/热度”代替；其二，\emph{不能}把“某方法在\emph{某项测试}上的结果”\emph{推广}为“任意序列、任意长度都有效”，\textbf{预印本里的观察只在它\emph{报告的具体设置}下成立}。

\section{知识演进的基本概念}
\label{sec:ket-overview-scenarios}
\label{sec:ket-overview-concepts}

先用同一个花卉系统，把“\textbf{变化到底发生在哪一层}”逐个看清——这是本章最重要的基本功，因为\textbf{不同层的变化，要改的对象和能保留的东西完全不同}。本章聚焦四类\textbf{典型的“随时间演进”场景}；至于“换一个全新任务类型”（如从认种跳到看病害）这种\textbf{跨任务迁移}，留给下一部分。

\textbf{场景一：环境随季节/地区漂移（\emph{观测分布}变了，但问题没变）}。系统最初用某地春季照片训练，运行半年后进入雨季、或搬到另一地区。注意这里\textbf{“要认的物种名称、判定标准可以\emph{完全不变}}”（还是认那十种山茶、茶梅），但\textbf{照片的\emph{分布}变了}：背景、光照（强弱/角度）、拍摄距离、各种花卉出现的比例都可能不同。于是旧模型里的线索\textbf{命运分化}：
\begin{itemize}\item \textbf{识别“\emph{花瓣边缘形状}}”这类\textbf{依赖物体本身的线索，或许仍可沿用}（花瓣形状不随季节变）；\item \textbf{依赖“\emph{当地背景颜色/光照}}”的线索（模型靠“这种花常衬在这种绿叶上、这种顺光下”来判断）则\textbf{可能\emph{失效}}。\end{itemize}
正确做法是\textbf{\emph{逐条检查“具体的预测依据}”还成不成立}，\textbf{避免两个相反的偷懒极端}：\textbf{既不能因为“任务名没变（还是认花）”就以为“无需任何调整”}，\textbf{也不能因为“外观看着变了”就“全部推倒重学”}——前者会带病部署、后者浪费已学到的物体特征。

\textbf{场景二：类别随时间\emph{陆续}增加（要学新类、还不能忘旧类）}。园区每隔一两个月引进一种新花卉。这和“一次性拿到所有类别再训练”\textbf{有本质区别}：新类别是\textbf{\emph{按顺序、随时间}到来的}，每来一批，系统都要在\emph{当时}可用的数据和资源下更新，并且\textbf{原来那十几种花还得继续认对}。这里立一条铁律：\textbf{“在新物种样本上\emph{答对了}”\emph{不能替代}“对\emph{原物种}重新测试}”——你\textbf{必须\emph{单独}重测旧任务}，新任务的高分不保证旧能力还在。神经网络在这种顺序更新下极易\textbf{“学会新的、覆盖旧的”}，这就是本章后面要反复讲的\textbf{灾难性遗忘}。

\textbf{场景三：观测设备更换（\emph{观测的含义}变了，连“输入坐标”都对不上）}。把普通 RGB 相机换成多光谱/红外传感器，\textbf{变化比“季节漂移”更根本}：多光谱每个通道测的是\textbf{不同波段}（可能含近红外等肉眼不可见的），和普通彩色 RGB 三通道\textbf{测的物理量完全不同}。这里有个\textbf{极易犯的错}：即使你把两边数据\textbf{都整理成“等长向量}”（硬凑成同样维度），\textbf{\emph{第$k$维}也\emph{不必然}表示“同一个观测属性}”——RGB 第一维是红光强度，多光谱第一维可能是近红外，\textbf{维度对齐\emph{并不等于}语义对齐}。所以即便\textbf{要输出的仍是“物种名称”（任务没变）}，也\textbf{必须为“新观测\emph{怎么接入}旧模型”\emph{建立依据}}；\textbf{\emph{单纯}“多训一阵子”\emph{不能}建立这种对应}——坐标语义的鸿沟不是梯度步数能自动填平的，需要显式的接入设计或对应证据。

\textbf{场景四：事实修正（\emph{连观测都没变}，是“关于世界的陈述”过期了）}。还有一类变化\textbf{根本不需要照片分布先改变}——\textbf{世界里的\emph{事实}本身更新了}。设系统附带一个植物研究机构的问答服务，机构负责人在某日由甲交接给乙：那么
\begin{itemize}\item \textbf{“\emph{现在}（交接后）问负责人是谁”}，应回答\textbf{乙}；\item \textbf{“交接\emph{之前}（历史上某时）负责人是谁”}，应回答\textbf{仍\emph{是甲}}。\end{itemize}
这类请求很特殊：\textbf{“要\emph{改变}什么、要\emph{保留}什么”其实\emph{更新请求本身已经规定好了}}（现任改、历史不改），难点在别处——\textbf{要决定“\emph{改哪个对象}}”：是改\textbf{外部事实条目}（知识库里的一条记录）、改\textbf{模型参数}、还是只改\textbf{当前这次问题的预测输入}（在提问里临时附上说明）？然后还要\textbf{\emph{按日期}检查回答对不对}。两种偷懒都不好：\textbf{完全沿用旧系统会留下\emph{过时答案}}（还在答甲）；\textbf{全部重训则可能\emph{重复大量无关工作}}（只是一个负责人字段变了，重训整个视觉模型毫无必要）。

图~\ref{fig:ket-overview-scenarios}把上面四类\textbf{随时间的变化}\textbf{画在\emph{同一个}已训练系统上}，\textbf{读法是“\emph{沿每条线看它\emph{改变的对象}是什么}”}。读图注意：\textbf{这些分支\emph{不是}“执行步骤”}，它们是\textbf{对“变化类型”的\emph{分类}}；而且\textbf{一次真实变化常常\emph{同时}踩多个分支}——比如“园区\emph{搬到新地区、那里还有没见过的花}”，就要\textbf{\emph{同时}处理“观测分布漂移”和“类别增量”}。\textbf{在\emph{选更新方法之前}，必须先\emph{说清}：哪些旧要求\emph{仍然有效}（要保留）、哪些关系\emph{需要重新建立}}。

\begin{figure*}[htbp]
  \centering
  % Conceptual relationships, not a temporal pipeline or measured gains.
% Final width is below 169 mm. All outgoing requests share the east-side port.
\begin{tikzpicture}[
  x=1mm,y=1mm,
  font=\figurefont\small,
  every node/.style={text=BookInk},
  source/.style={
    rectangle,draw=BookTeal,fill=BookTeal!8,
    line width=1pt,text width=39mm,minimum height=22mm,
    inner sep=2mm,align=center
  },
  request/.style={
    rectangle,draw=BookBlue,fill=BookPaper,
    line width=0.8pt,text width=55mm,minimum height=12mm,
    inner sep=1.5mm,align=center
  },
  link/.style={draw=BookTeal,line width=0.9pt},
  arrow/.style={link,-{Latex[length=2mm,width=1.4mm]}},
  change/.style={font=\figurefont\footnotesize,anchor=south,text=BookMuted}
]
  \node[source] (system) at (23,0) {持续运行的\\花卉识别系统};
  \node[font=\figurefont\footnotesize,align=center,text width=42mm]
    at (23,-18) {已有表示、模型与记忆\\附带机构问答服务};

  \node[request] (environment) at (124,27) {环境随季节/地区漂移\\光照、背景与类别比例变化};
  \node[request] (newclass) at (124,9) {类别陆续增加\\学新物种且不遗忘旧物种};
  \node[request] (sensor) at (124,-9) {观测设备更换\\RGB 换成多光谱或红外};
  \node[request] (editing) at (124,-27) {事实修正\\负责人交接与有效日期};

  \coordinate (fork) at (61,0);
  \draw[link] (system.east) -- (fork);
  \draw[link] (61,-27) -- (61,27);
  \foreach \target/\height/\object in {
    environment/27/观测分布,
    newclass/9/顺序学习与旧类保持,
    sensor/-9/观测含义与接入,
    editing/-27/事实有效范围
  } {
    \draw[arrow] (61,\height) -- (\target.west);
    \node[change] at (80,\height+1) {\object};
  }
  \node[font=\figurefont\footnotesize,text=BookMuted,anchor=north]
    at (95,-37) {各类变化可以同时发生};
\end{tikzpicture}

  \caption{已训练花卉系统随时间面对的不同变化。箭头指向使用请求，连线文字标明主要变化对象，不表示先后步骤。季节/地区漂移改变观测分布，类别陆续增加带来顺序学习与遗忘问题，传感器变化改变观测含义，事实修正对应附带问答服务中的负责人交接。一次真实变化可同时落在多个分支上。}
  \label{fig:ket-overview-scenarios}
\end{figure*}

还要防止“\textbf{一个名字掩盖多个条件}”：\textbf{同一次系统更新，可能\emph{同时}“换了地区”“只有少量新标签”“还要逐月持续学”}——这三句话\textbf{描述的是\emph{三件不同}的事}：\textbf{“换地区”说的是\emph{环境}（数据分布）；“少量标签”说的是\emph{监督信号}有多少；“逐月学习”说的是数据的\emph{到达顺序}}。如果你只用一个名字（比如笼统说“继续微调一下”）概括它，就\textbf{会\emph{漏掉}决定“某方法能不能用”的关键条件}。

所以接下来\textbf{8.1 用三件事把一个更新请求\emph{严格描述清楚}}：
\begin{enumerate}\item \textbf{先明确“已有的\emph{训练样本、模型、外部事实}”各是什么}（经验\emph{保存在哪}，8.1.1）；\item \textbf{再区分“\emph{实际改了哪个对象}”与“系统\emph{最终具备了什么能力}”}（8.1.1）；\item \textbf{最后区分“更新\emph{关系}”以及对“样本、模型、反馈”的\emph{访问权限}}（是回放、参数更新、还是外部条目修改，8.1.2–8.1.3）。\end{enumerate}

\subsection{知识与知识载体}
\label{sec:ket-overview-carriers}

要谈“维护一个已训练系统”，先得回答一个特别基础、却常被含糊带过的问题：\textbf{模型\emph{训练完}之后，所谓“学到的经验”到底\emph{以什么形式、存在哪}？}——因为\textbf{“更新时能动什么、遗忘会丢什么”完全取决于“你手里\emph{实际保留了什么}”}。训练一结束，\textbf{可供后续维护使用的对象至少有三类}，它们\textbf{形态不同、能回答的问题也不同}：
\begin{enumerate}\item \textbf{保留下来的\emph{训练样本及其标注}}（原始照片+物种标签，是“\emph{数据形态}”的经验）；\item \textbf{已经训练好的\emph{模型及其模型参数}}（一组数，是“\emph{压缩进模型}”的经验）；\item \textbf{可检索的\emph{事实条目或任务案例}}（外部知识库/记忆里“某机构负责人是乙”“某病斑长这样”等，是“\emph{显式记录}”的经验）。\end{enumerate}
\textbf{为什么非要分得这么细？}因为\textbf{它们\emph{不是}同一种东西，更新和遗忘的方式差别很大}：只有样本，你能回放重训但旧模型的现状不在其中；只有参数，模型能预测但你看不到原始数据、也指不出“哪个参数记着哪条事实”；只有事实条目，你能改“现任是谁”却改不动视觉识别。\textbf{本章始终\emph{区分}这些对象在演进中各承担什么作用}，\textbf{绝不含糊地统称成一种“\emph{没说清类型}的学习成果”}——类型不清，更新、遗忘和验证都无从谈起。

\textbf{先看“数据/记录形态”的两类}。\textbf{第一类，\emph{保留训练样本}}时，学习算法可以\textbf{\emph{重新读取}照片特征和标签}（想回放旧任务、纠错都有原料）；这正是对抗遗忘最直接的弹药。\textbf{第二类，使用\term{外部记忆}（external memory）}时，系统还能\textbf{检索物种描述、历史案例，以及\emph{带出处、时间、适用条件}的事实条目}（“某机构负责人是乙，自2026年7月1日生效”）——它像模型外挂的一本“可查阅、可增删改”的档案，区别于把一切都压进参数。

这里有两个\textbf{必须分清的界限}：\textbf{其一，“记忆里\emph{有}一份文档”\emph{不等于}“它成了\emph{训练样本}}”——\textbf{外部记忆中的文档\emph{不会自动}参与学习}，\textbf{只有当它\emph{实际进入训练目标的构造}时，才承担“训练样本”的职责}；否则它只是“可被检索到的资料”。\textbf{其二，不同载体\emph{信息不可互替}}：\textbf{只存一个物种名称，\emph{恢复不了}原照片的纹理}；\textbf{保存了事实条目，也\emph{不保证}检索器\emph{取到的是正确版本}}（可能取到过期的甲、而不是现任的乙）。所以要\textbf{\emph{分别}检查两件事}：\textbf{“条目\emph{有没有}正确保存（含生效日期）”}、以及\textbf{“这次预测\emph{有没有正确取用}它”}（存对了但取错，答案照样错）。

\textbf{再看“压缩进模型”的经验}。旧模型的参数是更新时最常被改动的对象，它有几种典型的被维护方式：可以\textbf{\emph{冻结}一部分、只训练新增加的部分}（学新类时不动底层，减少对旧能力的冲击）；可以\textbf{把旧参数当\emph{初始化}、允许继续训练}（从一个懂边缘纹理的位置继续优化，通常比随机初始化更快）；也可以\textbf{\emph{新增}任务分支或模块}（让新类别尽量走新参数，少覆盖旧参数）。这些方式的取舍（省不省、会不会忘、新类学不学得好）是第10章的主题，这里先建立“\textbf{参数是一种\emph{可被部分冻结、继续修改或旁路扩展}的载体}”这一直觉。

\begin{example}[带生效日期的事实修正]
\label{ex:ket-overview-dated-fact}
机构、人物和日期均为教学构造。设某植物研究机构在2026年7月1日由负责人甲交接给乙，此前一直由甲任职。8月查询现任负责人，预期答案是乙；查询2026年6月的负责人，预期答案仍是甲。这两个问题共同限定修改范围。

若更新外部事实条目，就应为甲与乙保存各自的有效日期，再检查问答是否取用了对应事实条目。若修改模型参数，就应检查不同问法以及交接前后的日期条件，模型参数文件改变本身不是回答正确的证据。若仅在当前问题前附上交接说明，变化只发生在这次预测可读取的预测输入中；除非另行保存或训练，下次独立查询没有理由自动取得这份说明。
\end{example}

这个带日期的例子逼出一对\textbf{必须分开的概念}：\textbf{“\emph{实际改了哪个对象}”} 和 \textbf{“系统最终\emph{具备了什么能力}”}。后者称为\term{任务能力}（task capability）——\textbf{指系统在\emph{规定的预测输入、可用工具、评价条件}下，\emph{实际能完成什么}}；而\textbf{实际修改对象}是\textbf{模型参数、外部事实条目、还是当前这次预测请求}（三者之一或组合）。

为什么要分？因为\textbf{“改对了对象”\emph{不蕴含}“能力对了”}，反之亦然：\textbf{就算检索\emph{已经返回了}“乙任职”的正确条目，仍要检查“最终回答\emph{有没有按日期\emph{正确使用}它}}”；反过来，\textbf{“有一次回答对了”也\emph{无法确定}这个对答案\emph{依据}来自哪}——是参数里记住的、检索文档命中的、还是这次提问临时附加的说明？\textbf{一次正确回答\emph{不能}反推依据来源}。正因如此，\textbf{第~\ref{chap:ket-update-maintenance}章会\emph{针对不同更新对象建立不同的验证}}：改参数要重测、改条目要测检索与取用、改上下文只在当次有效——\textbf{而\emph{不是}“看到文件变了/这次答对了”就验收通过}。

\textbf{还有一种最“抠门”的更新方式，\emph{连对象都不改}：只在当次预测的输入里临时附加说明}。把新物种的几张带答案示例\textbf{附在\emph{当前这次}提问里}（提示词中的少样本示例），模型当次可以照着预测。这些\textbf{当前\emph{可读取}的附加任务证据，称为\term{上下文}（context）}。\textbf{关键特性：\emph{撤去}示例后，\emph{下一次}预测\emph{没有义务}延续这种行为}——因为模型参数一点没改，它只是“\emph{这次看见了}”。所以要区分三个\textbf{改变位置}（持久化程度递增）：
\begin{itemize}\item \textbf{只改\emph{当前请求}（上下文）}：当次有效，下次无义务延续；\item \textbf{存进\emph{外部记忆}}：以后可再次检索读取，但要靠检索器取到；\item \textbf{训练、真正改动\emph{模型参数}}：\emph{永久}改变以后独立预测所走的计算。\end{itemize}
三者可结合，但\textbf{验证口径不同}：\textbf{“一次\emph{带着示例}答对了”，只证明了“\emph{当次被条件化}”时的能力}，要看是否真学会，必须\textbf{在\emph{撤去}示例的独立查询上重测}。

\subsection{知识演进的学习设定}
\label{sec:ket-overview-learning-settings}

把“变化如何随时间到来、系统要负什么责任”说清楚，就构成了演进问题的\textbf{学习设定}。先区分两种最容易被新手混为一谈的要求：

\textbf{其一是\emph{共同}学习（多个任务\emph{一起}训）}。如果物种、病害的数据\textbf{从一开始就都在手里}，可以让它们\textbf{共享同一套特征计算、损失一起回传}，这叫\term{多任务学习}（multi-task learning, MTL）。它\textbf{没有时间先后}，是“几个任务\emph{同时、对称}地互相塑造共享表示”。多任务的专门机制（共享哪些层、损失怎么平衡）属下一部分“知识的共享”，这里只需知道：\textbf{它研究的是“\emph{一起}学怎么协调”，\emph{不是}“后来的训练别冲掉前面的}”。

\textbf{其二是\emph{顺序}学习（任务/数据\emph{陆续}来）}，这才是本部分的核心。\textbf{当经验\emph{按顺序}不断到来}（先会认种、雨季漂移、再加新物种、再修事实，没完没了），学习者就必须\textbf{正面处理“\emph{新要求}与\emph{旧能力}的关系”}，这构成\term{持续学习}（continual learning，也叫终身学习）。用花卉系统说，持续学习要求系统在更新后同时做到：\textbf{新季节/新物种的新要求\emph{学会了}}，\textbf{旧物种的旧能力\emph{没丢}}，而且\textbf{更新若干轮之后\emph{仍然}学得动新东西}。

持续学习的\textbf{特殊之处在于\emph{时间}和\emph{有限资源}}：
\begin{itemize}\item \textbf{变化的\emph{来源}可以是环境漂移、类别增加、或（少见地）预测任务变化}；\item \textbf{“\emph{保留}哪些旧能力”必须\emph{显式指定}}：是只保核心物种、还是所有旧物种都要保？不能笼统说“别忘”，要落到具体旧任务的指标；\item \textbf{\emph{未来}的数据现在看不见}（不能拿秋天的数据帮春天做决策），\textbf{\emph{历史}数据还可能因存储/隐私\emph{被丢掉}}。\end{itemize}
\textbf{为什么“把所有历史数据存下来、每次\emph{完整重训}”\emph{不能当默认解}？}因为现实常常\textbf{两个条件\emph{不满足}}：\textbf{其一，\emph{历史照片\emph{不能全部保存}}}（隐私要求删除、存储有预算、数据在采集端带不走）；\textbf{其二，\emph{任务转换可能\emph{没有标记}}}（系统不知道“现在切换到新任务了”）。在这种“\emph{旧数据看不到/看不全、又要保旧能力}”的约束下，才逼出了持续学习的专门方法（第10章讲）。一个\textbf{必须排除的“假持续学习”}：\textbf{\emph{对同一批固定样本反复优化}（多训几个 epoch）只是\emph{求解过程}，\emph{训练步数变多}\emph{不构成}“持续学习”}——持续学习的关键是\textbf{\emph{有新的证据/要求\emph{随时间}进来}、且要\emph{维护}新旧能力}。

\textbf{持续学习里，“\emph{测试时给你多少信息}”会\emph{改变要解决的分类问题本身}}。以“类别陆续增加”为例，有三种\textbf{难度递增的标准设定}，区别全在“\emph{测试时\emph{告不告诉你}这是哪一批}”：
\begin{itemize}\item \term{任务增量学习}（task-incremental learning，Task-IL）：\textbf{测试时\emph{提供“任务身份}}”——明确告知“这是第 3 组物种里的一种”，于是你只需\textbf{在这一组内部\emph{对应的预测头}里选}。这是\textbf{最容易}的一档（帮你排除了其他组的类）。\item \term{类别增量学习}（class-incremental learning，Class-IL）：\textbf{类别逐步增加}（先认 10 种、下月再加 5 种），\textbf{测试时要在“\emph{截至当前的\emph{全部类别}}”里\emph{统一}选择}，\textbf{\emph{不给}任何能排除旧类的身份}。\textbf{难得多}：既要学新类、又不能和越来越多的旧类混淆，\textbf{遗忘问题最尖锐}。\item \term{领域增量学习}（domain-incremental learning，Domain-IL）：\textbf{类别和任务含义\emph{不变}}（始终认同样十种花），\textbf{只是\emph{依次经历不同的输入领域}}（春到夏再到秋、或不同地区/传感器），测试时\textbf{也不告诉你“该用哪个领域专属模型”}。\end{itemize}
\textbf{统一纪律}：\textbf{“任务身份、领域身份”在\emph{训练/选择/预测}三个环节\emph{各自可不可见}，必须\emph{逐项报告}}——\textbf{\emph{不能}用“我做的是持续学习”这种文献名\emph{代替}权限说明}；同一个算法在 Task-IL 上很好，换到 Class-IL 可能崩溃。

\subsection{知识演进的信息条件}
\label{sec:ket-overview-information}

\textbf{两个\emph{都叫}“持续更新”的方案，\emph{实际能拿到的样本、模型、反馈可能\emph{天差地别}}，能做的操作因此完全不同}——所以\textbf{比较更新方法前，必须先\emph{盘点访问权限}}。在持续维护里，最关键的是下面几类信息：
\begin{itemize}\item \textbf{\emph{历史样本}还能不能读到}：能，就可以\textbf{回放}旧任务；读不到（已删除/只有有界缓存），就只能靠参数约束或结构隔离防遗忘。\item \textbf{\emph{新数据}什么时候、以什么方式到}：是\textbf{离线攒成一批}（可反复遍历、统计全局），还是\textbf{在线逐张/逐阶段到达}（看不到未来、只能用已见部分）。\item \textbf{\emph{新标签}有多少、何时返回}：新物种的标注通常很少；鉴定结果若是“\emph{预测之后}才返回”（延迟反馈），那它\textbf{只能支持\emph{以后}的更新，\emph{不能}用于“\emph{此前}那次”预测}（你不能用明天才知道的答案改今天已发出的判断）。\item \textbf{\emph{模型参数}能不能改、能改多少}：可以全量微调，也可能因部署约束只能冻结主体、改一小部分（如适配器）或新增分支。\end{itemize}

\textbf{“\emph{要保留什么、有多少资源}”同样\emph{决定}哪些方案可行}。举两个\textbf{资源约束}的典型差别：\textbf{其一，\emph{总缓存固定}}（设备就这么大）时，每来一个新任务都要\textbf{\emph{淘汰}一部分旧内容}才能腾地方——回放哪些、忘哪些，本身就是核心决策；\textbf{其二，若“\emph{每个任务}给\emph{固定}缓存}”（可不断加存储），单任务不挤压，但\textbf{\emph{总存储会随任务数\emph{线性膨胀}}}，任务多了仍不可持续。\textbf{保留目标}也要说清：\textbf{“\emph{只}服务新季节”}（旧季节数据可弃）和\textbf{“新旧各季\emph{都}要继续服务”}（要防遗忘）\textbf{对应\emph{完全不同}的目标函数和方案}。

如果允许系统在“已经上线使用”阶段，根据当前到来的输入去改变自己的参数、统计量或其他保留状态，就进入\term{测试时适应}（test-time adaptation, TTA）。它的出发点很现实：部署后数据会持续漂移，与其等下次离线重训，不如让模型边用边调整（通常只用无标签输入、靠自监督损失更新）。但正因为发生在“测试/使用”阶段，TTA 协议必须把几件事写死：能不能看到完整批次（离线整批还是严格在线）、能不能跨样本保存状态（更新只对本张生效、还是累积到以后）、状态何时重置、有无反馈权限。\textbf{要特别区分“\emph{预测执行}（读当前这张照片出答案）”和“\emph{按目标样本调整模型}（更新参数或状态）}”：预测时当然要读当前输入，这并不等于被允许用一批新数据重新调参。

最后，\textbf{访问权限之外，还有\emph{授权、隐私、恶意来源}三道关}，须\textbf{分别}按\BookChapterRef{chap:trust-privacy}{卷3《基础、理论与可信性》第16章“隐私保护”}和\BookChapterRef{chap:trust-security}{卷3《基础、理论与可信性》第17章“对抗安全”}检查，\textbf{不能用“我做了访问限制”冒充“安全/隐私有保障”}。一个\textbf{反直觉但重要}的事实：\textbf{就算“训练样本\emph{不离开}采集地”或“只开放预测接口”，模型/接口\emph{仍可能泄露}训练样本的\emph{成员身份}或\emph{敏感属性}}（成员推断、属性推断攻击）。

\section{知识演进的研究内容}
\label{sec:ket-overview-research}

8.1 把“变化、载体、权限”描述清楚之后，本部分剩下三章就各司其职。下面用花卉系统把\textbf{演进的决策链}串起来，并给出章节地图。

\textbf{第一步：在\emph{真正更新前}，先判断“旧经验还值不值得用、用哪一部分”}。这件事由第~\ref{chap:ket-reuse-selection}章（第 9 章“知识的复用”）负责。它要回答：\textbf{旧模型里\emph{哪些层、哪些样本、哪些模块}在新季节/新类别上仍然有用}？是底层边缘特征通用、还是只是数据巧合？\textbf{值不值得复用，必须\textbf{相对于“\emph{不用}旧经验”的\emph{对照基线}}、在\emph{目标侧}衡量}——\textbf{旧季节 99\% 的准确率是\emph{来源侧}成绩，既不保证更新后仍然准、也不保证不会帮倒忙（负迁移）}。第9章还会处理“预算有限时先试训哪些候选”。这一章\textbf{既是演进的前置决策，也为下一部分的跨任务迁移提供“怎么判断可迁移性”的工具}。

\textbf{第二步：实施更新，让“\emph{新证据}与\emph{原有要求}\emph{同时}满足”}，由第~\ref{chap:ket-update-maintenance}章（第 10 章“知识的更新和维护”）负责。它要同时处理四件\emph{方向不同}的事：
\begin{itemize}\item \textbf{\emph{吸收}环境/观测变化}（雨季漂移、换传感器后让模型跟上新分布）；\item \textbf{\emph{保留}仍然有效的旧能力}（学会新物种的同时别忘了旧物种，即抗遗忘）；\item \textbf{\emph{修正}过时事实}（负责人甲换成乙，并保留历史日期的正确答案）；\item \textbf{按要求\emph{消除}某些训练样本的影响}（用户撤回授权后做“机器遗忘”）。\end{itemize}
这里要特别引入一个\textbf{只在“长期”尺度上才暴露的概念}——\term{可塑性}（plasticity）：\textbf{系统“\emph{继续}从\emph{后来}经验中\emph{学习}”的能力}。一个系统可能\textbf{“\emph{现在}旧任务都答对”，却已经\emph{僵化}、\emph{以后}很难再学新东西}（参数被压在饱和区、或被太强的防遗忘正则锁死），\textbf{当前答对\emph{不保证}以后仍易学}。所以长期维护要\textbf{把“保留旧能力（抗遗忘）”和“保持可学习性（可塑性）”\emph{分开}检查}（8.3 的可塑性损失会展开）。

还要\textbf{区分两个\emph{很像但验证标准完全不同}的操作}：\textbf{\emph{事实纠错}检验的是“新旧回答\emph{是否符合要求}}”（现任答乙、历史答甲）；\textbf{\emph{训练影响消除}则需要先\emph{指定删除标准}（到底要删到什么程度才算“该样本没影响过”），两者\emph{不能互换}}——“把输出里的甲改成乙”是事实编辑，\textbf{\emph{并不等于}“从训练影响意义上\emph{抹除}了甲相关样本”}。后者的验证\textbf{还要回到\emph{隐私保护}的专门边界}（成员推断、再训练近似等，见卷 3）。

\textbf{第三步：\emph{评价}一次更新到底成不成功}。这是持续学习评价里\textbf{最容易造假的一点}，务必分清两种测试口径：
\begin{itemize}\item \textbf{\emph{只在“本月新物种}”内部选择并计分}：这只检验\textbf{“\emph{新类之间}能不能区分”}，任务简单（候选少、还告诉你在新类里选）；\item \textbf{\emph{在“截至当前的\emph{旧+新所有}物种}共同集合”里选择并计分}：这才检验\textbf{“\emph{累积类别}上的分类”}（严格的类增量口径）。\end{itemize}
\textbf{危险在哪？}——\textbf{\emph{总体}错误率“达标”，仍可能\emph{掩盖旧类退化}}：模型可能在新类上表现好、把旧类大面积错判成新类，一平均却被新类的高分冲淡。\textbf{要判断旧能力是否\emph{真保留}，必须在“共同候选类别”口径下\emph{单列}旧类样本的错误率，并与更新\emph{前}的旧类错误率\emph{比较}}；比较时还要\textbf{保持三件事不变}：\textbf{旧类评价样本\emph{范围}不变、计分\emph{方式}不变、汇总\emph{权重}不变}，然后看\textbf{错误率\emph{增加量}是否超过\emph{预先规定的“保留容限}”}；\textbf{若要求“\emph{逐类}都保留”，还要\emph{分别}检查\emph{每一个}旧类}（不能只看旧类整体平均，某一种花可能已经全忘）。

这些跨阶段评价（能力保持矩阵、遗忘量、后向迁移、可塑性测试、事实编辑与机器遗忘的专门指标）集中在第~\ref{chap:ket-evolution-evaluation}章（第 11 章“知识演进的评价”）讲透。为了让你\textbf{遇到一个具体更新请求时能\emph{快速找到该读哪章}}，表~\ref{tab:ket-overview-chapters}\textbf{把本卷相关章节定位成一张地图}。

\begin{table*}[htbp]
  \centering
  \small
  \begin{tabularx}{\linewidth}{@{}p{52mm}XX@{}}
    \toprule
    章节 & 核心研究问题 & 主体节的组织依据 \\
    \midrule
    第~\ref{chap:knowledge-fusion}章：知识的融合 & 多来源学习依据与已有成果怎样共同使用 & 样本标注与预测输出、模块、模型参数、蒸馏及验证 \\
    第~\ref{chap:ket-overview}章：知识的演进 & 随时间变化提出什么问题，如何描述及定位研究 & 场景、载体与权限、研究问题与历史线索 \\
    第~\ref{chap:ket-reuse-selection}章：知识的复用 & 旧经验是否值得复用、哪一部分适用、应采用哪些 & 收益、适用性与候选选择 \\
    第~\ref{chap:ket-update-maintenance}章：知识的更新和维护 & 修改已有系统后，如何同时满足新旧要求 & 更新目标、证据、实施、验证与长期维护 \\
    第~\ref{chap:ket-evolution-evaluation}章：知识演进的评价 & 更新、保留、后续学习与指定修改是否满足要求 & 阶段指标、持续协议、事实修正及训练影响消除 \\
    第~\ref{chap:kts-overview}章：知识的迁移与共享概述 & 跨任务经验怎样帮助新任务、多任务怎样共训 & 基本概念、研究内容、研究历史 \\
    第~\ref{chap:ket-transfer-adaptation}章：知识的迁移 & 来源经验如何支持新目标及减少学习代价 & 目标确定、来源准备与目标适应 \\
    第~\ref{chap:ket-sharing}章：知识的共享 & 多个任务如何在共同学习中协调要求 & 任务组织、共享结构、联合目标与资源分配 \\
    第~\ref{chap:ket-evaluation}章：知识迁移与共享的评价 & 哪些证据说明迁移共享的预期收益已满足 & 指标、样本与基准、协议和证据分析 \\
    \bottomrule
  \end{tabularx}
  \caption{以直接研究目标组织的章节地图。中列给出要回答的问题，右列说明各章怎样分解问题；同一请求可同时涉及多章。第8–11章处理一个持续运行系统在时间维度上的演进，第12–15章处理跨任务的迁移与共享。}
  \label{tab:ket-overview-chapters}
\end{table*}

\textbf{这些章节\emph{不是}各管一段，它们之间存在\emph{真实的反馈回路}}：比如\textbf{\emph{为学新类而\emph{改动了共享底层}}，就必须\emph{重新测试}所有旧类}（共享部分一变，旧任务都受影响）；反过来，\textbf{\emph{完全冻结}底层虽然\emph{避免}覆盖旧计算，却可能\emph{限制}新类学习}（新类需要的细节学不进去）——这正是“学新”和“保旧”之间的核心张力。由此得到\textbf{评价的总原则}：\textbf{\emph{不能}只留下“某一个环节的最佳数值”}（比如只报告新类准确率新高），\textbf{必须\emph{同时}检查三样——\emph{新目标}达成没有、\emph{旧要求}保住没有、\emph{代价}花了多少}。

还有一个\textbf{会让查章变得困惑的事实}：\textbf{\emph{同一个“操作/技术名}”，\emph{不唯一}属于某一章}——它\textbf{服务于什么\emph{目的}，就归到相应的研究问题}。看四个在演进中常用的操作：
\begin{itemize}\item \term{知识蒸馏}（knowledge distillation）：用教师的预测或中间表示指导另一个模型。\textbf{目的不同、归属不同}——\textbf{用它\emph{压成一个承接模型}偏“融合”；用它\emph{把旧模型行为当作\emph{更新时的约束}}（让新模型别偏离旧模型太多以抗遗忘）就偏“更新维护”}。\item \term{经验回放}（experience replay）：让保存的历史样本\emph{再次参与训练}，让模型\emph{继续表达旧要求}（旧任务样本反复见，防止遗忘）——主要服务持续维护。\item \term{模型参数约束}（parameter constraint）：\textbf{限制新参数\emph{偏离旧参数}的程度}（如对重要参数加惩罚，8.3 的 EWC 即是），为“保留部分旧计算”提供正则约束。\item \term{模块路由}（module routing）：\textbf{按“当前输入属于哪个任务”\emph{选择}让哪些模块参与计算}（任务分支/混合专家），让不同任务尽量少覆盖同一些参数。\end{itemize}
但无论用了哪个操作，\textbf{都\emph{不能}因为“用了高级技术”就默认成功}——\textbf{采用任何操作，仍要\emph{分别验证}“\emph{新要求}学会没有、\emph{旧能力}保留好不好”}。

\textbf{现在给几个\emph{具体更新请求}“\emph{定主要职责}”}：
\begin{itemize}\item \textbf{\emph{撤回}一张训练照片（要求消除其影响）}：\textbf{主归“更新与维护”}，且要\textbf{套上\emph{可信性的删除标准}}（机器遗忘，见卷 3 隐私章）；\textbf{\emph{不能}只把缓存里的照片删掉就算数}——模型参数里可能仍留着它的影响，必须按删除标准验证。\item \textbf{\emph{雨季漂移后旧物种准确率下降}}：\textbf{主归“更新与维护”}（吸收分布变化 + 保住旧能力），先用第9章工具判断旧特征哪些仍有效。\item \textbf{\emph{每月加两种花卉、且保留旧类别}}：\textbf{在\emph{演进}口径下}属于持续的类别增量学习，评价要用前面讲的“单列旧类错误率 + 保留容限”；其中“旧知识怎么帮新类”的\emph{迁移手段}参考下一部分。\item \textbf{\emph{负责人由甲改任乙}}：\textbf{主归“更新与维护”的外部事实编辑}，按生效日期分别验证现任与历史。\end{itemize}
最后划清\textbf{本部分\emph{不}深入、但要知道边界}的两块：\textbf{其一，若“动作会\emph{改变后续状态分布}”、且决策好坏要看\emph{长期累计回报}}（比如主动选择下次标注什么、探索与利用的权衡），那已进入\textbf{序贯决策/强化学习}，可参考\BookChapterRef{chap:rl-overview}{卷7《学习范式：强化学习》第1章“强化学习概述”}；\textbf{其二，\emph{存储、通信、上线}的具体实现属于\emph{系统工程}}，本书不展开，但评估成本时要把它们算进去。

\section{知识演进的研究历史}
\label{sec:ket-overview-history}

“让一个系统随时间持续学习、又不忘记旧知识”这个愿望，几乎和联结主义本身一样老。历史上它主要沿\textbf{两条\emph{并行}、又互相纠缠的线索}发展：\textbf{其一是“\emph{顺序学习为何会忘、怎样不忘}”}（从灾难性遗忘到正则巩固、回放、结构隔离，再到可塑性）；\textbf{其二是“\emph{显式事实信念}怎样长期增删改”}（从永不停止的学习系统到知识编辑、机器遗忘）。读图/读史前先立三条规矩：其一，\textbf{线索\emph{长期并行}、互相交织}，不是先后接力；其二，\textbf{选某篇论文\emph{定位}某个思想，\emph{不等于}说它是该思想的“首次”}；其三，\textbf{\emph{后来}的方法\emph{并不}意味着\emph{替代}了前面全部方法}（EWC、回放这些“老”办法今天仍在用）。

\subsection{灾难性遗忘与抗遗忘机制}
\label{sec:ket-overview-history-maintenance}

回到场景二埋的伏笔：\textbf{新物种训练\emph{一旦改变了}“认旧物种”所依赖的那些参数，旧能力可能\emph{随之受损}}。这种现象有个专门名字——\term{灾难性遗忘}（catastrophic forgetting）：\textbf{用神经网络做\emph{顺序}学习时，\emph{后续}任务的训练可能使\emph{既有}能力\emph{大幅、甚至近乎完全}丧失}（不是缓慢衰减，而是被新梯度“覆盖”）。

两个要点：其一，\textbf{必须在\emph{原任务的同一评价口径}下检查}（同一批旧测试样本、同一种计分），\textbf{\emph{不能}凭“感觉”或“时间过了”判断}；其二，这\textbf{\emph{不是}大模型时代才有的问题}——French 在 1999 年的综述就\textbf{系统讨论了“表示重叠、学习干扰以及应对途径”}，说明\textbf{只要用共享参数顺序学多个任务，遗忘问题就存在}\scite{ket-french1999-forgetting}。

怎么对抗遗忘？历史上形成了\textbf{三大类思路}。\textbf{第一类是\emph{正则约束}}：一个自然想法是\textbf{“给\emph{重要}的参数上\emph{枷锁}”}，代表是\term{弹性权重巩固}（elastic weight consolidation, EWC），2017 年正式发表\scite{kirkpatrick2017ewc}。它先\textbf{估计“\emph{哪些参数对过去任务\emph{重要}}”}（用旧任务训练后这些参数附近的曲率/Fisher 信息近似），然后学新任务时：
\begin{itemize}\item 对\textbf{\emph{重要}参数偏离旧值，施加\emph{较大}代价}（像有根橡皮筋把它们拉回旧值，“弹性”由此得名）；\item 对\textbf{\emph{不重要}参数，保留\emph{较大}调整自由}。\end{itemize}
要准确理解它\textbf{保存的是什么}：\textbf{EWC 保存的是“\emph{历史参数取值} + \emph{参数重要性估计}”，\emph{不是}“\emph{全部}旧样本”}（不需要回放数据），\textbf{更\emph{不是}“某个参数\emph{一一对应}某条事实”}。本书取其 2017 年\emph{正式发表}年份（预印本 2016 年已公开）；且\textbf{原实验\emph{不保证}“任意长的任务序列都\emph{永不遗忘}”}——序列很长时，重要性估计会过时、枷锁也会累积。

\textbf{第二类是\emph{回放}（replay）}：把\textbf{旧样本、或\emph{替代}旧样本的\emph{生成}样本、或\emph{教师模型}在旧样本上的\emph{预测}，在学新任务时\emph{重新放回}训练}，让旧要求不断被重新表达。它\textbf{受“\emph{历史覆盖\emph{全不全}、缓存预算\emph{够不够}}”限制}（存不下全部旧样本时，回放什么、生成质量如何都成问题）。

\textbf{第三类是\emph{结构隔离}（structural / parameter isolation）}：\textbf{用“\emph{任务分支}”或“\emph{专属参数/模块}”，让不同任务\emph{尽量少\emph{直接覆盖}同一些参数}}（新任务走新分支，经典代表如渐进网络）。它\textbf{受“\emph{容量}（分支不能无限增）”和“\emph{模块\emph{怎么选}}（推理时走哪条分支）”条件限制}。

这三类方法逼出\textbf{一个\emph{权衡}和一对\emph{必须分开的检查}}：\textbf{参数惩罚（EWC）若\emph{过强}，会\emph{妨碍}新任务学习}（枷锁太紧，新东西学不进去）；回放占缓存、隔离占容量。所以\textbf{评价它们\emph{不能只看旧成绩}（那只反映遗忘）}，还要同时看：\textbf{其一，\emph{新要求}到底\emph{学会没有}（\textbf{可塑性}的当下表现）；其二，\emph{后续}任务\emph{是否仍能}在规定预算下学会（\textbf{可塑性}的长期表现）}。\textbf{这正是“\emph{遗忘检查}（旧的忘没忘）”和“\emph{可塑性检查}（新的还学不学得动）”的区别}——二者常常此消彼长，是持续学习几十年的核心张力。

\subsection{长期信念库、知识编辑与机器遗忘}
\label{sec:ket-overview-history-beliefs}

前面维护的还主要是“模型参数”，而一类持续运行系统把维护对象扩展到了“类别与关系构成的信念库”。Mitchell 等人在 2018 年的《Never-Ending Learning》中介绍了\term{永不停止的语言学习系统}（Never-Ending Language Learner, NELL）\scite{ket-mitchell2018-nell}。它日复一日地从网络文本中抽取“某实体属于某类别”“两实体间有某关系”，并为每条学到的“信念”保存置信度和出处，还能利用已学会的关系支持以后的抽取（旧信念约束新抽取，形成自我提升的循环）。时间口径要说清：论文报告系统自 2010 年起持续运行，2018 是“案例论文”发表年份。它对“长期维护”的概念贡献在于\textbf{拓宽了评价对象}：对这种系统，长期评价不能只看“某次训练后预测准不准”，还要看——实体与关系条目增长了多少、积累的错误怎么修正、学习过程本身有没有改进。

NELL 提醒我们信念会不断累积，但\textbf{事实条目也可能\emph{过时}，所以“保留”\emph{绝不等于}“一律冻结不变”}。这里出现两个\textbf{方向相反、却都要求“\emph{精确}改动”}的操作，务必分清：
\begin{itemize}\item \term{知识编辑}（knowledge editing）：\textbf{针对“\emph{指定的某条事实或某个行为}”做\emph{定向}修改}，同时\textbf{检验两件事——\emph{该改的改对了}（编辑成功）、\emph{不该改的\emph{其他}回答\emph{保持}住了}（locality，局部性）}。负责人交接例子正是如此：现任应答乙（要改），交接前历史应答甲（要保留），两者都要测。\item \term{机器遗忘}（machine unlearning）：\textbf{按\emph{预先指定的删除标准}，\emph{消除某些训练样本的影响}}。它的验收标准\emph{不是}“输出变了”，而是“\emph{影响}\emph{在规定意义上被消除}”。\end{itemize}
一个\textbf{关键澄清}：\textbf{“把回答里的甲改成乙”、或“以后\emph{不再返回}甲”，都\emph{没有直接验证}“甲相关样本的\emph{训练影响}已被\emph{消除}}”——前者是知识编辑能做到的表层改动，后者需要机器遗忘那套专门标准（影响分布是否逼近“从没训练过该样本”、能否抗成员推断等，属隐私验证）。

遗忘讲的是“旧的守不住”，长期维护还有\textbf{对称的另一极：“新的学不进”}。2024 年 Dohare 等人的研究\textbf{集中检验了}\term{可塑性损失}（loss of plasticity）：\textbf{系统在\emph{持续训练很久之后}，“\emph{学习\emph{新}要求}的能力”本身\emph{下降}了}（参数趋于饱和、激活退化、或被长期更新锁死）\scite{ket-dohare2024-plasticity}。他们用\textbf{长序列实验、并设置“\emph{重新初始化}网络”等对照}，得到几个警醒的发现：其一，\textbf{\emph{即使保留全部历史样本}（能一直回放），也\emph{不能自动}维持可塑性}；其二，\textbf{某些\emph{常用训练方法}在其所测设置里还会\emph{加重}可塑性下降}。论文提出的\textbf{持续反向传播（continual backpropagation）}做法是：\textbf{在常规梯度更新之外，\emph{逐步重置}那些“对当前计算\emph{贡献较小}”的单元}，腾出还能继续学习的新鲜单元；结果是\textbf{它在\emph{所测设置}中\emph{维持了可塑性}，但\emph{并没有同时解决}旧任务遗忘}——再次说明\textbf{“保旧”和“学新”是\emph{两个要分别评价}的目标}。结论\textbf{严格限定于所测网络、样本流和优化条件}，不能外推为万能方案。

这条线索走到这里，“模型参数、预测行为、显式列明的事实”被纳入同一个长期维护问题——它们都会随时间变、都要兼顾“改该改的、保该保的”，但\textbf{正因为载体不同，修改方式和验证标准不能共用一把尺子}：改参数要重测行为，改事实要测检索与取用，改输出口径并不代表消除了训练影响。其中实体关系图（知识图谱）的存储、本体（ontology）的版本管理等\textbf{显式知识的专门数据结构与工程}（图数据库、模式演化、一致性约束），本部分不深入，只讨论它们与“学习、检索、版本验证”的接口；修改范围怎么定、关联一致性怎么保、多次连续维护怎么不乱，在第~\ref{chap:ket-update-maintenance}章展开，“训练影响消除（机器遗忘）”仍须回到隐私保护的相应标准专门验证。

\subsection{顺序学习研究的两条共同边界}
\label{sec:ket-overview-history-lessons}

把上面两条历史线索合起来看，长期演进研究始终被两条共同边界约束，它们也直接对应第11章的评价设计：

\textbf{第一条边界：“学会新的”和“记住旧的”\emph{不能用同一个成绩互相替代}，而且还要再加一个“以后学不学得动”}。从 French 指出遗忘、EWC/回放/结构隔离三路上锁，到 Dohare 指出可塑性损失，历史反复证明：\textbf{任何只报告新任务成绩、或只报告旧任务保持的结论都是不完整的}；新类准确率、旧类错误率、后续新任务的学习曲线，是\textbf{三本要分别记的账}。

\textbf{第二条边界：改了“哪个载体”，就要用“匹配这个载体”的方式验证}。参数、外部事实、当次上下文的持久化程度不同：参数更新要看\emph{独立预测}是否改变，事实编辑要看\emph{检索与生效日期}是否正确，当次上下文\emph{撤去示例}后即失效，机器遗忘则要回到“\emph{影响是否消除}”的分布级标准。\textbf{“文件变了”“这次答对了”“名字不再出现了”，单独都不是验收通过的证据}。

图~\ref{fig:ket-overview-history}把本部分这两条线索的\textbf{代表工作}\textbf{画在\emph{同一根等距年份轴}上}。读图注意：横带\emph{只区分研究问题、不规定起止时间}，老方法至今仍在演化、被反复翻新；EWC 取 2017 年正式发表年份，NELL 取 2018 年案例论文年份。\textbf{至于“多任务共享规律”“快速适应/元学习”“多来源模型融合”这另外三条线索}（Caruana、Baxter、Ben-David、Yosinski、MAML、Hinton 蒸馏、Model Soups 等），它们主要回答的是\textbf{跨任务的迁移与共享}问题，已\textbf{移到下一部分第~\ref{chap:kts-overview}章的研究历史中展开}，本章不再重复。

\begin{figure*}[htbp]
  \centering
  % Equal year scale: x = 36 + 4 * (year - 1997), in mm.
% Event years are publication/case-paper years as noted in the text.
% Bands show concurrent research questions, not their starting/ending dates.
\begin{tikzpicture}[
  x=1mm,y=1mm,
  font=\figurefont\footnotesize,
  every node/.style={text=BookInk},
  band title/.style={
    align=center,text width=27mm,minimum height=10mm,
    inner sep=0.5mm,font=\figurefont\small
  },
  event/.style={align=center,inner sep=0.7mm},
  leader/.style={draw=BookMuted,line width=0.6pt}
]
  \foreach \level in {0,-32} {
    \fill[BookPaper] (33,\level-12) rectangle (162,\level+12);
    \draw[BookRule,line width=0.6pt] (35,\level) -- (158,\level);
  }
  \draw[BookMuted,line width=0.7pt] (36,19) -- (150,19);
  \foreach \year in {1997,2000,2005,2010,2015,2020,2025} {
    \draw[BookMuted,line width=0.6pt]
      ({36+4*(\year-1997)},18) -- ({36+4*(\year-1997)},20);
    \node[anchor=south] at ({36+4*(\year-1997)},20) {\year};
  }
  \node[anchor=east,text=BookMuted] at (29,21) {年份};

  \node[band title] (forgetting) at (14,0) {灾难性遗忘\\与抗遗忘机制};
  \node[band title] (beliefs) at (14,-32) {长期信念维护\\编辑与可塑性};

  \foreach \name/\year/\level/\ink in {
    french/1999/0/BookAmber,
    ewc/2017/0/BookAmber,
    replay/2017/0/BookAmber,
    nell/2018/-32/BookTeal,
    plasticity/2024/-32/BookTeal%
  } {
    \coordinate (\name) at ({36+4*(\year-1997)},\level);
    \fill[\ink] (\name) circle[radius=1mm];
  }

  \node[event,anchor=north] (french label) at (62,-4) {遗忘研究回顾\\French（1999）};
  \draw[leader] (french) -- (44.8,-2) -| (french label.north);
  \node[event,anchor=south] (ewc label) at (112,4) {EWC 正则巩固\\2017};
  \draw[leader] (ewc) -- (116,2) -| (ewc label.south);
  \node[event,anchor=north] (replay label) at (140,-4) {回放与结构隔离\\同期多路线};
  \draw[leader] (replay) -- (116,-2) -| (replay label.north);

  \node[event,anchor=north] (nell label) at (132,-36) {NELL 案例论文\\2018};
  \draw[leader] (nell) -- (120,-34) -| (nell label.north);
  \node[event,anchor=south] (plasticity label) at (146,-28) {知识编辑、机器遗忘\\与可塑性损失（2024）};
  \draw[leader] (plasticity) -- (144,-30) -| (plasticity label.south);
\end{tikzpicture}

  \caption{知识演进研究的两条并行线索中的代表工作。圆点按等距年份定位；横带区分“抗遗忘机制”与“长期信念维护”，不规定起止时间。EWC取2017年正式发表年份，NELL取2018年案例论文年份；图中位置不构成“首次”判断。跨任务的共享规律、快速适应与多来源融合见下一部分的概述章。}
  \label{fig:ket-overview-history}
\end{figure*}

\section{本章小结}
\label{sec:ket-overview-summary}

回到开头的花卉系统，那几个随时间到来的请求现在都有了各自明确的判断依据，而不再是笼统的“重训一下”：雨季漂移，要检查的是观测分布（照片怎么出现）与标签规律（给定照片后怎么定类）是否仍然支持旧预测，不能只凭“任务名没变”或“外观变了”下结论；逐月加入新物种，必须在“截至当前的累积类别”口径下同时评价新能力和旧能力，并单列旧类错误率对照保留容限；负责人交接说明：修改还可以由“事实的有效日期”触发，而不必等待照片分布发生任何变化，而且改参数、改外部条目、只改当次上下文是持久化程度完全不同的三件事。把这些请求归到哪类、归到哪一章，取决于“要实现什么、要保留什么、在什么权限和时间条件下做”，而不取决于“有没有再训练一次”——只改上下文、改事实条目也是正当的演进方式。

贯穿全章的一条方法论主线，是把“\textbf{载体、对象、证据、权限}”四件事分清楚。其一，先问经验保存在哪：保留下来的到底是训练样本、模型参数，还是可检索的外部事实条目；任务能力说的则是系统在规定输入、工具和评价条件下实际能完成什么——二者之间还隔着“怎么取用、怎么验证”的过程，改对对象不等于能力正确。其二，上下文使用、外部记忆保存、模型参数更新这三种改变持久化程度不同，要分别取证：当次条件化只在当次有效（撤去示例要重测），外部保存要验证以后能否检索取到正确版本，参数更新要验证独立预测是否改变。其三，也是演进区别于一般训练的根本：评价一次更新必须同时算三本账——新要求学会没有、旧能力保留多少、以后还学不学得动，三者不能用一个总分互相抵消。

历史上的两条线索——灾难性遗忘与抗遗忘机制（EWC 的正则约束、回放、结构隔离）、长期信念库与知识编辑/机器遗忘/NELL/可塑性损失——共同说明：让一个系统“持续变好又不垮掉”，从来不是单一技术能保证的，而是要在“学新”和“保旧”这对永恒张力之间，针对被修改的具体载体分别设计操作、分别建立验证。下一部分我们将离开“单一系统的时间线”，去看另一个互补的问题：一个任务上学到的经验，怎样跨过去帮助一个全新的任务（迁移），以及多个任务怎样从头就一起训练、共享同一套底子（共享）。这正是第~\ref{chap:kts-overview}章的内容。
